\documentclass[12pt]{article}
\usepackage[utf8]{inputenc}
\usepackage[T1]{fontenc}
\usepackage{hyperref}
\title{Yar\'i Language Resources for K-12 Curriculum Development}
\author{Compiled by Claude Sonnet 5.5 (claude-sonnet-5-5) for Andr\'as Kornai}
\date{}
\begin{document}
\maketitle

\section{Introduction}
Yar\'i (ISO 639-3 yri; Glottocode yari1235) is listed for the Yar\'i River area of Caquet\'a in the Colombian Amazon. The family label ``Bookkeeping'' in the survey metadata is a Glottolog placeholder, not a real language family; the language is effectively unclassified. Search results show it is poorly documented. Joshua Project reports roughly 800 people and says their language has not been classified, and some reports say they speak Spanish. An older Ethnologue-derived entry suggests it may be a dialect of Carijona (Carib), a Western Tucanoan language, or Huitoto. I could not verify any of these claims.

\section{Primary Educational Resources}
None found. No primer, textbook, school programme or policy document for Yar\'i appeared in search results.

\section{Linguistic and Reference Materials}
\begin{itemize}
\item Glottolog entry for Yar\'i (code and classification only): \url{https://glottolog.org/resource/languoid/id/yari1235.bigmap.html}
\end{itemize}
No grammar, wordlist or dictionary was found.

\section{Religious and Community Resources}
\begin{itemize}
\item Joshua Project, Yari in Colombia (population estimate, unclassified language): \url{https://joshuaproject.net/people_groups/22379}
\end{itemize}

\section{Digital and Audio Resources}
None found.

\section{Curriculum Development Context}
There is currently no evidence of any Yar\'i-language educational material, and the language's status (living, extinct, or a dialect of another language) is unclear from the sources found. Fieldwork would be needed before any curriculum could be planned. Two unrelated items appeared in the results and should not be confused with Yar\'i: a conlang-registry page for a ``Yari'' language, and the Carabayo language, which has been linked to Tikuna-Yur\'i.

\end{document}